% Pursuing Stacks (cotes 134-2 to 134-8): archive cote 134-5, site page 16 (PDF page 17), typescript page 203.
% typescript. Draft compared: ps4.tex:710-769.
% AI-transcribed by Claude Opus 5.5 (claude-opus-5-5, adaptive thinking, effort high), Message Batches API, batch ps2 (msgbatch_01NVwr1Yuo8qJf5rykeGBXPK), 2026-11-12; not yet checked by eye.
\scanpage{134-5/16}
\mspage{203}

(cohomological) \uline{smoothness} and (coh.) \uline{properness} of a map in (Cat), by
the obvious base-change propert\over{y}{ies}. It turned out that these could be readi-%
ly \over{ch}{cha}racterized by suitable asphericity conditions, which are formally
quite similar to the well known valuative criteri\over{um}{on} for a map of finite
type between noetherian schemes to be \deleted{either} ``universally open'' (which can
be viewed as a ``purely topological'' variant of the notion of smoothness),
\deleted{or}\inserted{resp.} ``universally closed'' \deleted{(for instance,} \inserted{or rather, more stringently,} proper\deleted{)}. These conditions, more-%
over, are trivially satisfied when f: X$\longrightarrow$Y turns X into a category over Y
which is \deleted{either \xout{co}}fibered over Y (i.e. definable in terms of a \inserted{contravariant} pseudo-func-%
tor $Y^{\circ}\longrightarrow$(Cat)) , \deleted{or} \inserted{resp.} \inserted{co}fibered over Y (namely, defin\xout{\unsure{a}}able in terms of a
\xout{contravariant} pseudo-functor $Y$\deleted{$^{\circ}$}$\longrightarrow$(Cat)). If we call such maps in (Cat)
``\xout{co}fibrations'' and ``\inserted{co}fibrations'' (very much in conflict, alas, with \xout{\unsure{the}}
Quillen's neat set-up of \xout{co}fibrations-\inserted{co}fibrations !) , it turns out that
\uline{fibrations are smooth}, \uline{cofibrations are proper}. This is all the more re-%
markab\unsure{k|l}e, as the notions of fibered and cofibered categories were intro-%
duced with a view upon ``large categories'', in order to pin down some
standard properties met in all situations of ``base change'' \hand{(}and \inserted{later} \hand{(}the dual
situation of co-base change \hand{)} - the main motivation for this being the
need to fomulate with a minimum of precision a set-up for \hand{``}descent techni-%
ques\hand{''} in algebraic geometry. (These techniques, as well as the cohomologi-%
cal base change theorems, make visibly a sense too in the context of
analytic spaces say, or of topological spaces, but they don't seem to
have been assimilated yet by geometers outside of algebraic geometry\hand{)}
That these typically ``general non-sense'' notions should have such precise
topological implications came as a complete surprise ! As a consequence,
a bifibration (namely a map which is both a fibration and a cofibration)
is smooth and proper and hence a (cohomological) ``Serre fibration'', for
instance the sheaves $R^{i}f_{*}(F)$, when F is a constant \inserted{abelian} sheaf on X \hand{(}i.e. a\xout{n}
\inserted{constant} abelian group object in $X^{\wedge}$\hand{)}, namely $y\hand{$\mapsto$}H^{i}(X_{/y},F)$ , are \uline{local} \uline{systems} \xout{\illegible{1}}
on Y i.e. factor through the fundamental groupoid of Y. \unsure{*}

A \xout{\unsure{still}} greater surprise still was the\xout{\illegible{1}} \uline{duality between the notions of
properness and smoothness}: just as a map f: X$\longrightarrow$Y in (Cat) is a cofibra-%
tion iff the ``dual'' map $f^{\circ}: X^{\circ}\longrightarrow Y^{\circ}$ is a fibration, it turns out that
f \uline{is proper iff} $f^{\circ}$ \uline{is smooth}. This was a really startling fact, and it
caused me to wonder, in the context of more general topoi than just those
of the type $X^{\wedge}$, whether there was'nt \uline{a notion of duality generalizing}
\uline{the relationship between two topoi $X^{\wedge}$ and} \hand{$X^{\vee}$} $\overset{\hand{def}}{=}$ $(X^{\circ})^{\wedge}$ . Indeed, the\unsure{se}
two categories of sheaves can be described intrinsically one in terms of
the other, up to equivalence, by a natural pairing

\margin{16}
